\subsection{Таймеры}
\begin{itemize}
  \item \verb|Timer.new(period, fn)| — создаёт \hyperlink{term:timer}{периодический таймер}; \verb|period| — \hyperlink{term:period}{период} в секундах; \verb|fn| вызывается каждый \hyperlink{term:tick}{тик}.
  \item \verb|timer:start()|, \verb|timer:stop()| — запуск/остановка таймера.
  \item \verb|Timer.callLater(delay, fn)| — единичный вызов \verb|fn| через \verb|delay| секунд.
  \item Таймеры и \hyperlink{term:state}{состояния} делайте \hyperlink{term:scope}{глобальными} (без \verb|local|), иначе сборщик мусора их удалит.
\end{itemize}
\textbf{Пример периодического таймера (мигание):}
\begin{codefile}{examples/chapter_05/leds_timers_blink.lua}
\end{codefile}
\textbf{Пример отложенного вызова (через 3 секунды):}
\begin{codefile}{examples/chapter_05/leds_timers_delayed.lua}
\end{codefile}
\textbf{Отложенный таймер внутри обработчика событий:}
\begin{codefile}{examples/chapter_05/leds_timers_callback_delayed.lua}
\end{codefile}
